% !TEX program = xelatex
% 来源与复核记录：
% Q1 改编自《基础30讲》高数第3讲导数定义例题链及1000题高数第3章第1题，检验复合自变量下的定义代换。
% Q2 改编自第3讲分段函数左右导数/可导条件例题链及1000题第3章第8题，答案用代数展开复核。
% Q3 改编自第3讲连续、左右导数与参数条件题型，答案由左右导数直接求解复核。
% Q4 改编自第3讲切线斜率与导数几何意义题型，答案由导数定义和切线公式复核。
\documentclass[11pt,a4paper,fontset=fandol]{ctexart}
\usepackage[a4paper,left=20mm,right=20mm,top=20mm,bottom=18mm,headheight=15pt]{geometry}
\usepackage{amsmath,amssymb,enumitem,fancyhdr,lastpage,needspace}
\setlength{\parindent}{0pt}\setlength{\parskip}{0.45em}\linespread{1.15}
\pagestyle{fancy}\fancyhf{}\fancyhead[L]{\small 11408 / 22408}\fancyhead[R]{\small 数学二}\fancyfoot[C]{\small 第 \thepage\ 页，共 \pageref{LastPage} 页}\renewcommand{\headrulewidth}{0.6pt}
\newcommand{\workarea}[1]{\par\vspace{0.2em}\noindent\rule{\linewidth}{0.35pt}\vspace{#1}\par}
\begin{document}
\section*{第 3 讲收尾检测}
\begin{enumerate}[label=\textbf{\arabic*.},itemsep=0.9em]
\item 设函数 $f$ 在 $0$ 处可导，且 $f(0)=0$。求
\[
\lim_{x\to0}\frac{f(\sin x)}{x}.
\]
\workarea{5.0cm}

\item 设
\[
f(x)=\begin{cases}
ax^2+bx+1,&x\le 1,\\
\ln x+c(x-1),&x>1.
\end{cases}
\]
若 $f$ 在 $x=1$ 处可导，求 $a,b,c$ 之间应满足的全部条件。
\workarea{6.0cm}

\item 设
\[
f(x)=\begin{cases}
\dfrac{\sin(ax)}{x},&x\ne0,\\
b,&x=0,
\end{cases}
\]
其中 $a,b$ 为常数。求使 $f$ 在 $x=0$ 处可导的 $a,b$。
\workarea{5.5cm}

\item 曲线 $y=f(x)$ 在点 $(1,2)$ 处有切线。若
\[
\lim_{h\to0}\frac{f(1+2h)-f(1)}{h}=6,
\]
求该切线方程。
\workarea{5.0cm}
\end{enumerate}
\vfill
\begin{center}\small 第 3 讲：一元函数微分学的概念\end{center}
\end{document}
